% !TeX root = ../../../main.tex
기룡이는 이번에 열린 알고리즘 대회에 관심을 가지게 됐다.
대회를 주최한 동아리의 이름은 SOLUTIO(솔루티오)이다.
AI 컴퓨터공학부 학생들이 모여서 만든 알고리즘 동아리다.
SOLUTIO가 마음에 들었던 기룡이는 SOLUTIO와 관련된 프로그램을 하나 만들어보려고 한다.

자연수 $n$을 입력받아, $n$이 짝수라면 ``\texttt{KGU SOLUTIO}''를 출력하고,
$n$이 홀수라면 ``\texttt{KGU AI CSE}''를 출력하는 간단한 프로그램이다.
참고로 AI CSE는 ``AI 컴퓨터공학부''의 영문명인 ``AI Computer Science and Engineering''의 약자이다.

아직 프로그래밍 초보자인 기룡이를 도와 원하는 프로그램을 작성해 보자.

참고: 아래 사진은 솔루티오의 마스코트 솔리(Solly)이다. 솔리의 부리를 쓰다듬으면 문제가 잘 풀린다는 소문이 있다.

\begin{center}
    \includegraphics[width=8.0cm]{SOLUTIO 로고.png}
\end{center}
